\paperchapter{October 2011}{October 2011}{A}{60}{GSET 2011}
\keytable{\kcell{1}{C} & \kcell{2}{D} & \kcell{3}{C} & \kcell{4}{A} & \kcell{5}{A} & \kcell{6}{C} & \kcell{7}{C} & \kcell{8}{B} & \kcell{9}{B} & \kcell{10}{C} \\
\kcell{11}{D} & \kcell{12}{C} & \kcell{13}{A} & \kcell{14}{C} & \kcell{15}{A} & \kcell{16}{A} & \kcell{17}{D} & \kcell{18}{B} & \kcell{19}{C} & \kcell{20}{C} \\
\kcell{21}{C} & \kcell{22}{C} & \kcell{23}{A} & \kcell{24}{A} & \kcell{25}{D} & \kcell{26}{C} & \kcell{27}{A} & \kcell{28}{D} & \kcell{29}{C} & \kcell{30}{B} \\
\kcell{31}{B} & \kcell{32}{C} & \kcell{33}{D} & \kcell{34}{C} & \kcell{35}{C} & \kcell{36}{C} & \kcell{37}{C} & \kcell{38}{B} & \kcell{39}{D} & \kcell{40}{A} \\
\kcell{41}{B} & \kcell{42}{B} & \kcell{43}{A} & \kcell{44}{B} & \kcell{45}{C} & \kcell{46}{C} & \kcell{47}{B} & \kcell{48}{A} & \kcell{49}{B} & \kcell{50}{B} \\
\kcell{51}{A} & \kcell{52}{B} & \kcell{53}{D} & \kcell{54}{A} & \kcell{55}{C} & \kcell{56}{A} & \kcell{57}{C} & \kcell{58}{C} & \kcell{59}{B} & \kcell{60}{D}}
\begin{context}
\textbf{Read the following text and answer the Question Nos.~1 to 6.}\par\smallskip
The algorithmic theory of beauty takes the subjectivity of the observer into account and postulates that among several observations classified as comparable by a given subjective observer, the aesthetically most pleasing one is the one with the shortest description. This is closely related to the principles of algorithmic information theory and minimum description length. For example, mathematicians welcome simple proof with a short description in their formal language. Another concrete example describes an aesthetically pleasing human face whose proportions can be described by very few bits of information, drawing inspiration from 15th century proportion studies by Leonardo da Vinci. Going further, this theory explicitly distinguishes between what is beautiful and what is interesting. Interestingness corresponds to the first derivative of subjectively perceived beauty. Here, the premise is that any observer continually tries to improve the predictability of the observations by discovering regularities such as repetition and symmetry.
\end{context}
\begin{question}{1}
{According to the algorithmic theory:\par}
\optsone{subjective perceptions of observers' cannot be systematically classified}{there is no correlation between subjective observation and beauty}{there is a positive correlation between beauty and minimum description}{perception of beauty is unpredictable}
\answer{C}{there is a positive correlation between beauty and minimum description}
\begin{explain}
The theory says the most pleasing observation is the one with the \textbf{shortest description}. Beauty therefore increases as description length decreases: there is a \textbf{positive correlation between beauty and minimum description}.
\end{explain}
\end{question}

\begin{question}{2}
{The formal language of mathematics\par}
\optsone{prefers simple proof which may or may not be elaborate}{prefers the least elaborate description and repetition}{may be complex even as it provides a short description}{ensures simplicity and bravity}
\answer{D}{ensures simplicity and bravity}
\begin{explain}
Mathematicians ``welcome simple proof with a short description in their formal language''. The formal language of mathematics therefore values \textbf{simplicity and brevity}.
\end{explain}
\end{question}

\begin{question}{3}
{The author of the passage suggests that Leonardo da Vinci:\par}
\optsone{was a precursor of the algorithmic theory of information}{predicted the algorithmic theory of beauty}{produced proportion studies that support algorithmic theory of beauty}{gave a theory of subjectively perceived aesthetics}
\answer{C}{produced proportion studies that support algorithmic theory of beauty}
\begin{explain}
The passage says the face example draws ``inspiration from 15th century \textbf{proportion studies} by Leonardo da Vinci''. His studies \emph{support} the theory; he did not predict it.
\end{explain}
\end{question}

\begin{question}{4}
{Which of the following is not true ?\par}
\optsone{What is interesting is not always beautiful}{We perceive beauty when we discover symmetry}{Mathematicians prefer short descriptions}{We can distinguish between what is interesting and what is beautiful}
\answer{A}{What is interesting is not always beautiful}
\begin{explain}
The official key marks (A) as the statement not supported. The passage says only that the theory \emph{distinguishes} the interesting from the beautiful, with interestingness being the \emph{rate of change} (first derivative) of perceived beauty. It does not claim that what is interesting is not beautiful.
\end{explain}
\end{question}

\begin{question}{5}
{How do we perceive beauty ?\par}
\optstwo{By noticing patterns of similarity}{By subjective observation}{By identifying a pleasing human face}{By producing the shortest description}
\answer{A}{By noticing patterns of similarity}
\begin{explain}
The observer improves predictability ``by discovering regularities such as \textbf{repetition and symmetry}'', i.e.\ by noticing patterns of similarity.
\end{explain}
\end{question}

\begin{question}{6}
{An appropriate title for the passage would be :\par}
\optstwo{Leonardo da Vinci on Beauty}{Aesthetics and Mathematics}{Algorithmic theory of Beauty}{The Beautiful and the Interesting}
\answer{C}{Algorithmic theory of Beauty}
\begin{explain}
The passage is entirely about the \textbf{algorithmic theory of beauty}.
\end{explain}
\end{question}

\begin{question}{7}
{The gas leaked during Bhopal disaster was:\par}
\optstwo{ethyl cyanoacetate}{ethyl cyanoacrylate}{methyl isocyanate}{methyl diisocyanate}
\answer{C}{methyl isocyanate}
\begin{explain}
The Bhopal gas tragedy (2--3 December 1984, Union Carbide plant) released \textbf{methyl isocyanate (MIC)}, killing thousands. It is the world's worst industrial disaster.
\end{explain}
\end{question}

\begin{question}{8}
{The audible range of frequency for human ear is\par}
\optstwo{10 Hz to 10 kHz}{20 Hz to 20 kHz}{10 Hz to 20 kHz}{10 Hz to 30 kHz}
\answer{B}{20 Hz to 20 kHz}
\begin{explain}
The human audible range is about \textbf{20 Hz to 20 kHz}. Below it is infrasound; above it is ultrasound.
\end{explain}
\end{question}

\begin{question}{9}
{If waste materials contaminate the source of drinking water, which of the following diseases will spread?\par}
\optsfour{Scurvy}{Typhoid}{Malaria}{Anaemia}
\answer{B}{Typhoid}
\begin{explain}
Contaminated drinking water spreads water-borne diseases such as \textbf{typhoid}, cholera and hepatitis A.
\why{Scurvy is vitamin-C deficiency; malaria is mosquito-borne; anaemia is a nutritional deficiency.}
\end{explain}
\end{question}

\begin{question}{10}
{Which of the following is most likely to limit primary production in fresh water lakes?\par}
\optsfour{Oxygen}{Calcium}{Phosphorus}{Carbon}
\answer{C}{Phosphorus}
\begin{explain}
In fresh-water lakes \textbf{phosphorus} is usually the limiting nutrient for algal (primary) production. Excess phosphate from detergents and fertilisers causes eutrophication.
\end{explain}
\end{question}

\begin{question}{11}
{Which of the following is not a major greenhouse gas?\par}
\optstwo{Carbon dioxide}{Methane}{Water vapour}{Calcium carbonate}
\answer{D}{Calcium carbonate}
\begin{explain}
CO$_2$, methane and water vapour are greenhouse gases. \textbf{Calcium carbonate} is a solid (limestone), not a gas.
\end{explain}
\end{question}

\begin{question}{12}
{Tsunami is caused due to :\par}
\optsone{Giant high tide waves in ocean}{Oceanic bad weather conditions}{Sudden motion of the ocean floor}{Intercontinental collisions in the ocean}
\answer{C}{Sudden motion of the ocean floor}
\begin{explain}
A tsunami is caused by the \textbf{sudden displacement (motion) of the ocean floor}, usually by an undersea earthquake, landslide or volcanic eruption.
\end{explain}
\end{question}

\begin{question}{13}
{The full form of FTP is:\par}
\optstwo{File Transfer Protocol}{File Transfer Program}{File Transfer Property}{File Transfer Problem}
\answer{A}{File Transfer Protocol}
\begin{explain}
\textbf{FTP} is \textbf{File Transfer Protocol}, used to transfer files between computers over a network.
\end{explain}
\end{question}

\begin{question}{14}
{For proper functioning of computers, which of the following is not recommended ?\par}
\optsone{Anti-virus should always be activated}{Always use authorized softwares}{Never connect to internet}{Never allow external data devices on your PC without checking for viruses}
\answer{C}{Never connect to internet}
\begin{explain}
Anti-virus, authorised software and scanning external devices are recommended. \textbf{``Never connect to the internet''} is not a recommendation; the internet is essential, and safe use is what matters.
\end{explain}
\end{question}

\begin{question}{15}
{If '10110110' is written on a paper, what does it mean?\par}
\optsfour{Data}{Information}{Knowledge}{Wisdom}
\answer{A}{Data}
\begin{explain}
A bare string of bits such as 10110110, without context or processing, is \textbf{data}.
\tip{DIKW hierarchy: Data $\rightarrow$ Information $\rightarrow$ Knowledge $\rightarrow$ Wisdom.}
\end{explain}
\end{question}

\begin{question}{16}
{Which of the following is not a characteristic of information ?\par}
\optstwo{It is unchangeable with time}{It is meaningful to the receiver}{It is relevant to the activity}{It is clear and concise}
\answer{A}{It is unchangeable with time}
\begin{explain}
Information is meaningful, relevant, clear and \textbf{timely}. Its value \emph{changes with time}, so ``unchangeable with time'' is not a characteristic of information.
\end{explain}
\end{question}

\begin{question}{17}
{The selection of information and communication technological tools should essentially depend on:\par}
\optstwo{Technological infrastructure}{Learner's access to the technology}{The support personnel}{All of the above}
\answer{D}{All of the above}
\begin{explain}
The choice of ICT tools depends on the available infrastructure, learners' access and support staff, i.e.\ \textbf{all of the above}.
\end{explain}
\end{question}

\begin{question}{18}
{Which of the following is a soft copy output device ?\par}
\optstwo{Printer}{Visual Display Unit (VDU)}{CD-ROM}{Both (A) and (B) of the above}
\answer{B}{Visual Display Unit (VDU)}
\begin{explain}
A \textbf{soft copy} is displayed electronically, not printed. The \textbf{VDU (monitor)} is a soft-copy output device; a printer gives hard copy; a CD-ROM is a storage medium.
\end{explain}
\end{question}

\begin{question}{19}
{Which of the following is correct sequence of steps in scientific method?\par}
\optsone{Observation, Deduction of consequences, Hypothesizing, Verification}{Verification, Hypothesizing, Observation, Deduction of consequences}{Observation, Hypothesizing, Deduction of consequences, Verification}{Hypothesizing, Observation, Verification, Deduction of consequences}
\answer{C}{Observation, Hypothesizing, Deduction of consequences, Verification}
\begin{explain}
Scientific method: \textbf{observation $\rightarrow$ hypothesis $\rightarrow$ deduction of consequences $\rightarrow$ verification}.
\end{explain}
\end{question}

\begin{question}{20}
{A research carried out by a practitioner is called:\par}
\optstwo{Applied Research}{Basic Research}{Action Research}{Experimental Research}
\answer{C}{Action Research}
\begin{explain}
Research by a practitioner (teacher, manager) to improve his or her own practice is \textbf{action research}.
\end{explain}
\end{question}

\begin{question}{21}
{Sampling procedure in which every subject of a population has equal chance of being selected is called:\par}
\optstwo{Purposive sampling}{Cluster sampling}{Probability sampling}{Non-probability sampling}
\answer{C}{Probability sampling}
\begin{explain}
When every unit has an equal (known) chance of selection, the procedure is \textbf{probability sampling}.
\end{explain}
\end{question}

\begin{question}{22}
{In research, reliability indicates:\par}
\optstwo{Precision in measurement}{Error in measurement}{Replicability in measurement}{Generalisability in measurement}
\answer{C}{Replicability in measurement}
\begin{explain}
\textbf{Reliability} is the consistency of measurement: the same results on \textbf{replication} (repeated measurement).
\end{explain}
\end{question}

\begin{question}{23}
{An example of non-empirical research is :\par}
\optsfour{Philosophical}{Historical}{Survey}{Experimental}
\answer{A}{Philosophical}
\begin{explain}
\textbf{Philosophical research} relies on reasoning and conceptual analysis rather than empirical data, so it is non-empirical.
\end{explain}
\end{question}

\begin{question}{24}
{Delimitation in research report refers to:\par}
\optsone{Setting limits to the problem of research}{Limitation of the researcher in studying the problem}{Defining the research problem}{Alternative to scope of the problem}
\answer{A}{Setting limits to the problem of research}
\begin{explain}
\textbf{Delimitations} are the boundaries the researcher deliberately sets on the study (population, variables, area, period), i.e.\ \textbf{setting limits to the problem}. \emph{Limitations} are constraints beyond the researcher's control.
\end{explain}
\end{question}

\begin{question}{25}
{The main purpose of classroom teaching is to:\par}
\optsone{prepare students for examination}{make students content expert}{develop competitive spirit among students}{make students apply disciplinary knowledge to life}
\answer{D}{make students apply disciplinary knowledge to life}
\begin{explain}
Classroom teaching aims to help students \textbf{apply disciplinary knowledge to life}, not merely to pass examinations or compete.
\end{explain}
\end{question}

\begin{question}{26}
{Which of the following has potential to influence learning ?\par}
\optstwo{Socio-economic status}{Ability to express}{Study habits}{Ability of books}
\answer{C}{Study habits}
\begin{explain}
Per the key, \textbf{study habits} are the learner factor with the potential to influence learning directly.
\end{explain}
\end{question}

\begin{question}{27}
{Which of the following factors contributes towards effective teaching ?\par}
\optsone{Knowledge of teaching methods and learners}{Knowledge of content and institution}{Knowledge of teaching method and institution}{Knowledge of learner and computer}
\answer{A}{Knowledge of teaching methods and learners}
\begin{explain}
Effective teaching requires \textbf{knowledge of teaching methods and of learners} (pedagogy and psychology), along with content.
\end{explain}
\end{question}

\begin{question}{28}
{Which of the following is the advantage of lecture method ?\par}
\optstwo{Effective face-to-face teaching}{Ensures attention of all students}{Caters to all individual differences}{Suitable for large group teaching}
\answer{D}{Suitable for large group teaching}
\begin{explain}
The lecture method's main advantage is that it is \textbf{suitable for teaching large groups} economically.
\end{explain}
\end{question}

\begin{question}{29}
{Teachers use teaching aids, because:\par}
\optsone{Some teachers are incapable of teaching on their own}{Students require some change while learning}{It is necessary to reach the students' level of understanding}{To use the updated technology of teaching}
\answer{C}{It is necessary to reach the students' level of understanding}
\begin{explain}
Teaching aids are used to reach \textbf{the students' level of understanding}, making ideas concrete and clear.
\end{explain}
\end{question}

\begin{question}{30}
{Internal assessment marks are to be:\par}
\optstwo{kept confidential}{revealed to students}{based on attendance}{used as gracing marks}
\answer{B}{revealed to students}
\begin{explain}
Internal assessment marks should be \textbf{revealed to students} so that they get feedback and the process stays transparent.
\end{explain}
\end{question}

\begin{question}{31}
{Efficacy and effectiveness of communication can be assessed through:\par}
\optsfour{Reception}{Feedback}{Transmission}{Encoding}
\answer{B}{Feedback}
\begin{explain}
The effectiveness of communication is judged by \textbf{feedback}, which shows whether the message was understood as intended.
\end{explain}
\end{question}

\begin{question}{32}
{Which one of the following aspects has no influence in face-to-face communication ?\par}
\optstwo{Language ability}{Facial expressions}{Location}{Acquaintance}
\answer{C}{Location}
\begin{explain}
Language, facial expressions and acquaintance all influence face-to-face communication. The key treats \textbf{location} as having no influence.
\end{explain}
\end{question}

\begin{question}{33}
{Which of the following is not a type of communication ?\par}
\optstwo{Vertical communication}{Horizontal communication}{Crosswise communication}{Verbal communication}
\answer{D}{Verbal communication}
\begin{explain}
Vertical, horizontal and crosswise (diagonal) are the \emph{directions} of communication flow in an organisation. The key gives \textbf{verbal communication} as the answer: it is classified by \emph{mode}, not by direction, so it does not belong to this set of types.
\end{explain}
\end{question}

\begin{question}{34}
{Noise in communication is to:\par}
\optstwo{Block the media}{Distort the message codes}{Add on the message codes}{Distort the decoding}
\answer{C}{Add on the message codes}
\begin{explain}
Per the key, noise \textbf{adds on to the message codes}: unwanted signals get added to the message and distort it.
\end{explain}
\end{question}

\begin{question}{35}
{Which of the following is a pre-requisite for effective classroom communication ?\par}
\optstwo{Reflective questioning}{Elaborate explanation}{Drawing attention of students}{Planning the lesson}
\answer{C}{Drawing attention of students}
\begin{explain}
The first requirement for classroom communication is \textbf{drawing the students' attention}. Without attention no message is received.
\end{explain}
\end{question}

\begin{question}{36}
{Which of the following is not a feedback loop in classroom communication?\par}
\optstwo{Answering a question}{Taking down the notes}{Creating disturbance}{Cross-questioning the teacher}
\answer{C}{Creating disturbance}
\begin{explain}
Answering questions, taking notes and cross-questioning the teacher are responses that feed back to the teacher. \textbf{Creating disturbance} is not a feedback loop in classroom communication.
\end{explain}
\end{question}

\begin{question}{37}
{Participation of students representatives in higher education governing bodies is advisable to:\par}
\optsone{treat students community at par with teacher}{prepare future politicians}{promote democratic citizenship}{provide extra privilege to subside student unrest}
\answer{C}{promote democratic citizenship}
\begin{explain}
Student representation in governing bodies helps \textbf{promote democratic citizenship}: participation, responsibility and a voice in decisions.
\end{explain}
\end{question}

\begin{question}{38}
{Researches in the area of environmental education fall under:\par}
\optstwo{Inter-disciplinary research}{Multi-disciplinary research}{Intra-disciplinary research}{Uni-disciplinary research}
\answer{B}{Multi-disciplinary research}
\begin{explain}
Environmental education draws on biology, chemistry, geography, economics and sociology, so it is \textbf{multi-disciplinary} research.
\end{explain}
\end{question}

\begin{question}{39}
{The purpose of distance education in India is to:\par}
\optstwo{provide degrees economically}{generate large number of graduates}{increase the rate of higher education}{reach the unreachables}
\answer{D}{reach the unreachables}
\begin{explain}
Distance education aims to \textbf{reach the unreached}: people who cannot attend regular institutions because of distance, work or circumstances.
\end{explain}
\end{question}

\begin{question}{40}
{Teaching at higher education is:\par}
\optstwo{not a job, though salaried}{not a time based activity}{not a contractual activity}{not a cost effective activity}
\answer{A}{not a job, though salaried}
\begin{explain}
Teaching in higher education is a vocation and mission: it is \textbf{not merely a job, though salaried}.
\end{explain}
\end{question}

\begin{question}{41}
{How can teacher inculcate values among students ?\par}
\optsone{By imposing discipline}{By behaving ideally}{By telling life histories of great people}{By telling moral stories}
\answer{B}{By behaving ideally}
\begin{explain}
Values are best taught by example: the teacher must \textbf{behave ideally} (``values are caught, not taught'').
\end{explain}
\end{question}

\begin{question}{42}
{Educational institutes' accrediting bodies are:\par}
\optsone{Government bodies to control and regulate Government policies}{Statutory bodies to maintain quality of education}{Bodies to control management indulging in large profits}{Bodies to improve examination system of educational institution}
\answer{B}{Statutory bodies to maintain quality of education}
\begin{explain}
Accrediting bodies (NAAC, NBA) are \textbf{statutory or autonomous bodies to maintain and assure the quality of education}.
\end{explain}
\end{question}

\begin{question}{43}
{Complete the sequence of squares by choosing the correct figure from the given alternatives:\\[3pt]
\centering
\newcommand{\sq}[1]{\framebox[1.1cm][c]{\rule{0pt}{0.8cm}\parbox[c]{1cm}{\centering\tiny #1}}}
\sq{$\angle$}\,\sq{$\triangle\angle$}\,\sq{$\triangle\triangle\triangle$}\,\sq{$\triangle\triangle\triangle$\\$\triangle\angle$}\,\sq{$\triangle\triangle\triangle$\\$\triangle\triangle\triangle$\\$\angle$}\,\sq{?}\par}
\optsfour{\framebox{\parbox{0.9cm}{\centering\tiny$\triangle\triangle\triangle$\\$\triangle\triangle\triangle$\\$\triangle\triangle\triangle$}}}{\framebox{\parbox{0.9cm}{\centering\tiny$\triangle\triangle\triangle$\\$\triangle\triangle\triangle$\\$\triangle\triangle\triangle$\\$\triangle\angle$}}}{\framebox{\parbox{0.9cm}{\centering\tiny$\triangle\triangle\triangle$\\$\triangle\triangle\triangle$\\$\triangle\triangle\triangle$\\$\triangle\triangle\angle$}}}{\framebox{\parbox{0.9cm}{\centering\tiny$\triangle\triangle\triangle$\\$\triangle\triangle\triangle$\\$\triangle\angle\angle$}}}
\answer{A}{\framebox{\parbox{0.9cm}{\centering\tiny$\triangle\triangle\triangle$\\$\triangle\triangle\triangle$\\$\triangle\triangle\triangle$}}}
\begin{explain}
The number of complete triangles grows $0, 1, 3, 4, 6, \ldots$ and the figure alternates between having a half-drawn triangle ($\angle$) and not having one. In the 3rd figure the angle is ``completed'' into a full row of triangles; likewise the next step completes the figure into \textbf{three full rows, i.e.\ nine triangles}, as in option (A).
\end{explain}
\end{question}

\begin{question}{44}
{Put the correct letter for `?' in the following diagram:\\[3pt]
\centering\begin{tabular}{ccc}
\framebox{C} & $\triangleright$ B & \textcircled{F}\\[2pt]
\framebox{D} & $\triangleright$ C & \textcircled{?}\\[2pt]
\framebox{E} & $\triangleright$ D & \textcircled{T}\\
\end{tabular}\par}
\optsfour{H}{L}{M}{P}
\answer{B}{L}
\begin{explain}
The circle letter's position is the \textbf{product} of the other two positions:
\[\text{C}(3)\times\text{B}(2)=6=\text{F},\qquad \text{E}(5)\times\text{D}(4)=20=\text{T}\]
\[\text{D}(4)\times\text{C}(3)=12=\mathbf{L}\]
\end{explain}
\end{question}

\begin{question}{45}
{Take Universal set U $=\{$A, B, C, \ldots, Z$\}$ of all English alphabets. If set A contains the letters from GUJARAT, the members of set B are the alphabets from PUNJAB, the set C contains the alphabets from SIKKIM and the elements of the set D are the alphabets from TAMILNADU, then the Venn diagram for the sets A, B, C, D will be:\par}
\optsfour{\tikz[baseline=-0.5ex,scale=0.25]{\draw (-2,-2) rectangle (2,2); \draw (0,0.8) circle(0.8); \draw (-0.8,0) circle(0.8); \draw (0.8,0) circle(0.8); \draw (0,-0.8) circle(0.8);}}{\tikz[baseline=-0.5ex,scale=0.25]{\draw (-2,-2) rectangle (2,2); \draw (-0.6,1.1) circle(0.7); \draw (-0.6,0) circle(0.7); \draw (-0.6,-1.1) circle(0.7); \draw (0.7,0.1) circle(0.7);}}{\tikz[baseline=-0.5ex,scale=0.25]{\draw (-2,-2) rectangle (2,2); \draw (-0.9,-0.5) circle(0.7); \draw (0.1,0) circle(0.7); \draw (0.8,0.6) circle(0.7); \draw (0.7,-0.2) circle(0.7);}}{\tikz[baseline=-0.5ex,scale=0.25]{\draw (-2,-2) rectangle (2,2); \draw (0,0.8) circle(0.9); \draw (-0.2,1.0) circle(0.35); \draw (0,-0.8) circle(0.9); \draw (-0.2,-1.1) circle(0.35);}}
\answer{C}{\tikz[baseline=-0.5ex,scale=0.25]{\draw (-2,-2) rectangle (2,2); \draw (-0.9,-0.5) circle(0.7); \draw (0.1,0) circle(0.7); \draw (0.8,0.6) circle(0.7); \draw (0.7,-0.2) circle(0.7);}}
\begin{explain}
Write out the sets:
\begin{center}\small
\begin{tabular}{ll}
A (GUJARAT) & \{G, U, J, A, R, T\}\\ B (PUNJAB) & \{P, U, N, J, A, B\}\\ C (SIKKIM) & \{S, I, K, M\}\\ D (TAMILNADU) & \{T, A, M, I, L, N, D, U\}\\
\end{tabular}
\end{center}
A$\cap$B $=$ \{U, J, A\}, A$\cap$D $=$ \{T, A, U\}, B$\cap$D $=$ \{U, N, A\}, and A$\cap$B$\cap$D $=$ \{U, A\}. C shares \{I, M\} with D only and nothing with A or B.
So A, B, D overlap mutually and C overlaps only D, as in diagram \textbf{(C)}.
\end{explain}
\end{question}

\begin{question}{46}
{A fighter pilot notices a row of several tanks of enemy. He attacks on the tank T at the middle in the row. He then attacks on one more tank T1 in that row immediately preceding to T. In order to finish the last tank in that row, he has to attack 9th tank counting in order from T1. The total number of tanks in that row is:\par}
\optsfour{13}{11}{17}{15}
\answer{C}{17}
\begin{explain}
T is the middle tank and T1 is the tank immediately before it. From T1 the pilot counts 9 more tanks to reach the last one, so there are $9-1=8$ tanks after T. Since T is in the middle, there are also 8 tanks before T.
Total $=8+1+8=\mathbf{17}$.
\end{explain}
\end{question}

\begin{question}{47}
{Fill in the cell of `?' by putting the correct letter:\\[3pt]
\centering\begin{tabular}{|c|c|c|}\hline G & B & C\\\hline H & F & H\\\hline A & D & ?\\\hline\end{tabular}\par}
\optsfour{G}{E}{P}{N}
\answer{B}{E}
\begin{explain}
Per the key the answer is \textbf{E}. With letter positions, the first and third columns then have the same total, $G+H+A=7+8+1=16=C+H+E=3+8+5$, balancing about the middle column.
\end{explain}
\end{question}

\begin{question}{48}
{Which member of the following set differs from the other members ? (OX, GIRAFFE, CAMEL, DOG, GOAT, ELEPHANT, DEER)\par}
\optsfour{DOG}{DEER}{CAMEL}{OX}
\answer{A}{DOG}
\begin{explain}
Ox, giraffe, camel, goat, elephant and deer are \textbf{herbivores}; the \textbf{dog} is a carnivore (and a pet), so it is the odd one out.
\end{explain}
\end{question}

\begin{question}{49}
{Put the number in the empty box by choosing the correct digit from the given alternatives:\\[3pt]
\centering\begin{tabular}{|c|c|c|}\hline 7 & 3 & 6\\\hline 5 & 4 & 7\\\hline 8 & ? & 2\\\hline\end{tabular}\par}
\optsfour{9}{6}{5}{3}
\answer{B}{6}
\begin{explain}
Each row adds up to 16: $7+3+6=16$ and $5+4+7=16$, so $8+?+2=16$ gives $?=\mathbf{6}$.
\end{explain}
\end{question}

\begin{question}{50}
{If $a*b=\dfrac{b}{a}$, then which of the following statements is \emph{not} true? ($a$ and $b$ are real numbers)\par}
\optsfour{$a*0=0,\ a\neq0$}{$0*b=0,\ b\neq0$}{$1*1=1$}{$a*a=1,\ a\neq0$}
\answer{B}{$0*b=0,\ b\neq0$}
\begin{explain}
$0*b=\dfrac{b}{0}$ is \textbf{undefined} (division by zero), so (B) is not true.
Check the others: $a*0=0/a=0$ \checkmark; $1*1=1$ \checkmark; $a*a=a/a=1$ \checkmark.
\end{explain}
\end{question}

\begin{question}{51}
{For which set of values of $p, q, P, Q$ and for which arithmetic operation $*$, the expression $\left(\dfrac{p}{P}\right)*\left(\dfrac{q}{Q}\right)=-1$?\par}
\optstwo{p = -1, q = 5, P = Q = -4 and * is '+'}{p = P = 1, q = Q = -1 and * is 'x'}{p = Q = -1, q = P = 1 and * is '-'}{p = q = P = 1, Q = -1 and * is '/'}
\answer{A}{p = -1, q = 5, P = Q = -4 and * is '+'}
\begin{explain}
Option (A): $\dfrac{-1}{-4}+\dfrac{5}{-4}=\dfrac14-\dfrac54=-1$ \checkmark.
(B) gives $1\times1=1$; (C) gives $-1\div-1=1$; (D) gives $1-(-1)=2$.
\end{explain}
\end{question}

\begin{question}{52}
{Complete the sequence by putting the appropriate number from the given alternatives: 12, 15, 16, 20, 20, 25, 24, ?\par}
\optsfour{24}{30}{27}{35}
\answer{B}{30}
\begin{explain}
Two interleaved series: $12, 16, 20, 24$ ($+4$) and $15, 20, 25, \mathbf{30}$ ($+5$).
\end{explain}
\end{question}

\begin{question}{53}
{What is the appropriate missing letter in the following sequence ? A, B, D, H, ?\par}
\optsfour{I}{L}{M}{P}
\answer{D}{P}
\begin{explain}
Positions double: A(1), B(2), D(4), H(8), \textbf{P(16)}.
\end{explain}
\end{question}

\begin{question}{54}
{If $l=\sqrt3$, $m=1/7$, $n=2$ and $t=5/4$, then $l^2 * m \;\square\; n \;o\; t=-11/28$ for\par}
\optstwo{$*$, $\square$ and $o$ are all `$-$'}{$*$ is $+$, $\square$ is $+$, $o$ is `$-$'}{$*$ is $-$, $\square$ is $-$, $o$ is `$+$'}{$*$, $\square$ and $o$ are all `$+$'}
\answer{A}{$*$, $\square$ and $o$ are all `$-$'}
\begin{explain}
With all three operations as minus: $3-\dfrac17-2-\dfrac54=\dfrac{84-4-56-35}{28}=-\dfrac{11}{28}$ \checkmark.
\end{explain}
\end{question}

\begin{context}
Answer the questions 55, 56 and 57 from the given data in graphical form. Given data of number of washing machines in thousands, manufactured by two companies A and B, during 1997 to 2002.
\begin{center}
\begin{tikzpicture}
\begin{axis}[width=10cm,height=5cm,xmin=1996.5,xmax=2002.5,ymin=40,ymax=175,xtick={1997,...,2002},
  x tick label style={/pgf/number format/1000 sep=},tick label style={font=\scriptsize},ylabel={\scriptsize thousands},
  nodes near coords,every node near coord/.append style={font=\tiny},legend style={font=\tiny,at={(1.02,1)},anchor=north west}]
\addplot[mark=square,thick,brand] coordinates {(1997,139) (1998,120) (1999,100) (2000,128) (2001,107) (2002,148)};
\addplot[mark=o,thick,dashed,accent] coordinates {(1997,119) (1998,99) (1999,141) (2000,78) (2001,120) (2002,159)};
\legend{A,B}
\end{axis}
\end{tikzpicture}
\end{center}
\end{context}
\begin{question}{55}
{What is the difference between the number of washing machines manufactured by company B in 1997 and 1998?\par}
\optsfour{10,000}{11,000}{20,000}{30,000}
\answer{C}{20,000}
\begin{explain}
B: $119-99=20$ thousand $=\mathbf{20{,}000}$.
\end{explain}
\end{question}

\begin{question}{56}
{What is approximate average number of washing machines manufactured by company A over the given period?\par}
\optsfour{123666}{741999}{431555}{310111}
\answer{A}{123666}
\begin{explain}
A: $\dfrac{139+120+100+128+107+148}{6}=\dfrac{742}{6}\approx123.67$ thousand $\approx\mathbf{123666}$.
\end{explain}
\end{question}

\begin{question}{57}
{In which year the difference between the production of companies A and B was minimum?\par}
\optsfour{1997}{1998}{2002}{2001}
\answer{C}{2002}
\begin{explain}
Differences (thousand): 1997: 20; 1998: 21; 1999: 41; 2000: 50; 2001: 13; 2002: \textbf{11}. The minimum is in \textbf{2002}.
\end{explain}
\end{question}

\begin{context}
Given the table of marks obtained by two candidates P and Q in four papers. Answer the following questions 58, 59 and 60 from this data.
\begin{center}\small
\begin{tabular}{lcccc}\toprule \textbf{Paper} & I & II & III & IV\\\midrule
Maximum marks & 80 & 70 & 100 & 60\\ Candidate P & 48 & 56 & 75 & 42\\ Candidate Q & 40 & 63 & 82 & 36\\\bottomrule
\end{tabular}
\end{center}
\end{context}
\begin{question}{58}
{Arrange in increasing order the \% marks obtained by the candidate P. The sequence of papers is:\par}
\optsfour{I, II, IV, III}{I, III, II, IV}{I, IV, III, II}{II, III, I, IV}
\answer{C}{I, IV, III, II}
\begin{explain}
P's percentages: I $=48/80=60\%$, II $=56/70=80\%$, III $=75\%$, IV $=42/60=70\%$.
Increasing order: \textbf{I (60), IV (70), III (75), II (80)}.
\end{explain}
\end{question}

\begin{question}{59}
{In which paper(s) the candidate Q obtained the distinction?\par}
\optsfour{III}{II, III}{II, III, IV}{II}
\answer{B}{II, III}
\begin{explain}
Q's percentages: I $=50\%$, II $=63/70=90\%$, III $=82\%$, IV $=60\%$. Distinction ($\ge75\%$) in \textbf{II and III}.
\end{explain}
\end{question}

\begin{question}{60}
{Approximate \% results of candidates P and Q will respectively, be:\par}
\optsfour{70\% and 80\%}{70\% and 60\%}{70\% and 50\%}{71\% and 70\%}
\answer{D}{71\% and 70\%}
\begin{explain}
Both candidates total $221/310$: P $=48+56+75+42=221$ and Q $=40+63+82+36=221$, so each scores about $71.3\%$. Among the options, the key takes \textbf{71\% and 70\%} as the approximate results.
\end{explain}
\end{question}
